% -------------------------------------------------------------------------
% WBS ID: RES-PHY-BAR
% Deliverable: Barrier–Flow Interaction and Energy Dissipation
% Status: In progress
% Evidence status: Local barrier-flow output roles established
% Review status: Not reviewed
% Last updated: 2026-07-19
% Dependencies: RES-PHY-SCL, RES-PHY-LOD, RES-MOD-EQS
% Downstream interfaces: RES-NUM-SPEC, RES-OPT-OBJ, RES-ML-REP, RES-VER-MET
% -------------------------------------------------------------------------

\section{RES-PHY-BAR --- Barrier–Flow Interaction and Energy Dissipation}
\label{sec:res-phy-bar}

\subsection{Purpose and Scope}
\label{sec:res-phy-bar-purpose}

% TODO: Define what this Deliverable covers and explicitly excludes.

This Deliverable links local barrier-flow mechanisms to the physical
outputs extracted from the local three-dimensional URANS--VOF model.
It defines what is meant by comparative hydrodynamic barrier
performance before structural failure or final certification are
considered. The active geometry scope is fixed-rigid comparison of
wall-type and obstacle/dissipating/array-type barriers.

\subsection{Research Questions}
\label{sec:res-phy-bar-research-questions}

% TODO: State the questions that must be answered before the Deliverable
% can be accepted.

\subsection{Terminology and Definitions}
\label{sec:res-phy-bar-definitions}

% TODO: Define the technical terminology, symbols, quantities and
% classifications used in this Deliverable.

\subsection{Literature Evidence}
\label{sec:res-phy-bar-literature-evidence}

% TODO: Record evidence from peer-reviewed papers, authoritative datasets,
% standards, technical reports and primary documentation.
%
% Do not create a detached literature-summary list. Explain how each source
% supports, contradicts or limits a project decision.

\textcite{kozelkov2022turbulence} supports the claim that turbulence
can materially affect run-up, overtopping and barrier loading in
Navier--Stokes tsunami simulations. \textcite{deljesus2012porous}
supports deferring unresolved porous-barrier physics to a separate
activation-gated model. Source roles: `3D-SST`, `3D-RANSVOF` and
`3D-DEFERRED`. Project conclusion: the baseline comparison shall use
resolved hydrodynamic outputs for impermeable or explicitly resolved
barrier geometries, while unresolved porous dissipation is not claimed
unless a porous model is activated.

\subsection{Governing Relations and Quantitative Definitions}
\label{sec:res-phy-bar-governing-relations}

% TODO: Record relevant equations, dimensionless groups, empirical models,
% extraction rules or quantitative definitions.
%
% Use align environments for displayed equations.
% Every displayed equation must have a unique \label{eq:...}.
% Do not place punctuation after displayed equations.

Overtopping volume during a loading interval is:

\begin{align}
  V_{\mathrm{ot}}
  &=
  \int_{t_0}^{t_1}
  \int_{\Gamma_{\mathrm{crest}}}
  \max
  \left(
    \mathbf{u}\cdot\mathbf{n}_{\mathrm{crest}},\,
    0
  \right)
  \alpha
  \,\dd\Gamma
  \,\dd t
  \label{eq:res-phy-bar-overtopping-volume}
\end{align}

A transmission measure may be defined from incident and transmitted
energy-like wave measures:

\begin{align}
  K_{\mathrm{T}}
  &=
  \frac{
    E_{\mathrm{trans}}
  }{
    E_{\mathrm{inc}}
    +
    \epsilon_E
  }
  \label{eq:res-phy-bar-transmission-measure}
\end{align}

Reflection and apparent dissipation measures shall be reported with
the same incident-energy normalisation:

\begin{align}
  K_{\mathrm{R}}
  &=
  \frac{
    E_{\mathrm{refl}}
  }{
    E_{\mathrm{inc}}
    +
    \epsilon_E
  }
  \label{eq:res-phy-bar-reflection-measure}
  \\
  D_E
  &=
  1
  -
  K_{\mathrm{T}}
  -
  K_{\mathrm{R}}
  -
  \frac{
    E_{\mathrm{ot}}
  }{
    E_{\mathrm{inc}}
    +
    \epsilon_E
  }
  \label{eq:res-phy-bar-apparent-dissipation-measure}
\end{align}

Decision role: these equations define comparison diagnostics for
energy redirection and loss, not a closed mechanical energy proof.
Supporting evidence: \textcite{deljesus2012porous} for barrier
transmission/reflection distinctions and \textcite{kozelkov2022turbulence}
for turbulence-sensitive run-up and overtopping. Limitation: final
energy definitions must be compatible with available validation
measurements and numerical outputs.

\subsection{Physical Mechanisms}
\label{sec:res-phy-bar-physical-mechanisms}

% TODO: Complete this RES-PHY Domain-specific subsection.

The barrier-flow mechanisms of interest are incident-wave run-up,
reflection, overtopping, flow separation, wake formation, local
recirculation, transmitted wave reduction, pressure loading and
turbulent dissipation. For the baseline impermeable model, energy
dissipation is inferred from resolved hydrodynamic changes and
turbulence modelling rather than unresolved porous resistance.

Wall-type barriers primarily test reflection, run-up, overtopping,
crest overflow and loading. Obstacle, dissipating and array barriers
primarily test blockage, gap jets, separation, wake interaction,
transmission reduction and downstream momentum redistribution. Source
role: this distinction connects RES-GEO classes to measurable flow
mechanisms. Supporting citations: \textcite{deljesus2012porous} and
\textcite{watanabe2022validation}. Downstream handoff: RES-GEO-MET and
RES-VER-MET shall preserve class-specific energy and loading metrics.

\subsection{Characteristic Spatial and Temporal Scales}
\label{sec:res-phy-bar-characteristic-spatial-temporal-scales}

% TODO: Complete this RES-PHY Domain-specific subsection.

\subsection{Relevant Observables}
\label{sec:res-phy-bar-relevant-observables}

% TODO: Complete this RES-PHY Domain-specific subsection.

Relevant observables are maximum run-up, overtopping volume or
discharge, transmitted wave height or energy measure, reflected-wave
measure, downstream maximum depth and velocity, peak pressure, force,
impulse, moment and wake/recirculation indicators.

\subsection{Controlling Dimensionless Parameters}
\label{sec:res-phy-bar-controlling-dimensionless-parameters}

% TODO: Complete this RES-PHY Domain-specific subsection.

\subsection{Physical Regimes and Transition Conditions}
\label{sec:res-phy-bar-physical-regimes-transition-conditions}

% TODO: Complete this RES-PHY Domain-specific subsection.

\subsection{Implications for Model Validity}
\label{sec:res-phy-bar-model-validity-implications}

% TODO: Complete this RES-PHY Domain-specific subsection.

\subsection{Candidate Approaches}
\label{sec:res-phy-bar-candidate-approaches}

% TODO: Define the credible candidate approaches considered.

\subsection{Critical Comparison}
\label{sec:res-phy-bar-critical-comparison}

% TODO: Compare candidates using physically and project-relevant criteria.
% Do not select an approach solely on popularity or implementation ease.

\subsection{Assumptions and Applicability}
\label{sec:res-phy-bar-assumptions}

% TODO: State assumptions and the conditions under which they are valid.

\subsection{Limitations and Failure Conditions}
\label{sec:res-phy-bar-limitations}

% TODO: State where the evidence, model or selected approach ceases to be
% reliable or applicable.

\subsection{Project-Specific Conclusions}
\label{sec:res-phy-bar-conclusions}

% TODO: Convert the evidence into conclusions specific to the Studentship
% project. Do not merely restate literature findings.

\subsection{Selected Approach and Rationale}
\label{sec:res-phy-bar-selected-approach}

% TODO: Record the selected approach only after sufficient evidence exists.
% State the alternatives rejected and the reasons for rejection.

The selected barrier-performance approach is comparative and
hydrodynamic: candidate defences are compared using run-up,
overtopping, transmission, reflected energy, pressure, force, impulse,
moment and downstream-flow metrics extracted from accepted local 3D
simulations. The word accepted here refers to timestep/run acceptance
and later verification status; this Deliverable does not claim that
validation has been executed.

\subsection{Unresolved Decisions}
\label{sec:res-phy-bar-unresolved-decisions}

% TODO: Record unresolved questions, required evidence and decision owners.

Open decisions are the final energy diagnostic, normalisation of
transmission/reflection measures, accepted overtopping threshold,
porous-extension activation gate and weighting of hydrodynamic metrics
for optimisation.

\subsection{Requirements and Handoffs}
\label{sec:res-phy-bar-requirements}

% TODO: State requirements passed to other Research Domains, Software,
% verification activities or later execution work.

`RES-NUM-SPEC` shall extract these metrics after timestep acceptance.
`RES-OPT-OBJ` shall decide how hydrodynamic outputs become objective
functions. `RES-ML-REP` shall define which metrics become surrogate
outputs, and `RES-VER-MET` shall define validation and acceptance
criteria.

\subsection{Deliverable Completion Record}
\label{sec:res-phy-bar-completion-record}

% TODO: Record whether the Deliverable is:
%
% - Not started
% - In progress
% - Draft complete
% - Reviewed
% - Accepted
%
% Acceptance requires:
%
% - the research questions to be answered;
% - claims to be supported by appropriate evidence;
% - assumptions and limitations to be explicit;
% - the project-specific conclusion to be stated;
% - downstream requirements to be traceable;
% - unresolved decisions to be closed or formally deferred.
